\chapter{Arquitectura del Sistema y Trazabilidad}
\label{ch:arquitectura}

Este capítulo documenta la arquitectura de software y el mecanismo de trazabilidad que vincula cada resultado cuantitativo con su artefacto computacional de origen.

\section{Arquitectura Modular}

El sistema se organiza en cinco módulos con interfaces explícitas, lo que permite probar cada módulo de forma aislada y ejecutar experimentos en paralelo sin interferencia.

El módulo de \textbf{datos} expone una interfaz común para los tres conjuntos. El módulo de \textbf{simulación de daños} retorna la máscara binaria junto con la semilla que la generó, garantizando reproducibilidad determinista. El módulo de \textbf{interpolación} contiene las dieciocho implementaciones bajo una clase base abstracta con interfaz unificada, eliminando ventajas espurias por acceso diferencial a metadatos. El módulo de \textbf{evaluación} calcula las seis métricas. El módulo de \textbf{optimización} orquesta Optuna con registro estructurado de métricas intermedias, puntos de control y detección temprana de convergencia.

\section{Cadena de Trazabilidad}

Cada ejecución sigue una cadena determinista: carga de configuración, carga de datos, simulación de daño con semilla fija, interpolación, evaluación con seis métricas, y registro estructurado de resultados con metadatos completos. Cada ejecución recibe un identificador único (SHA-256 de la tupla completa de parámetros) que permite verificar reproducibilidad bit a bit, rastrear cualquier número reportado hasta su archivo de origen, y detectar contaminación entre fases experimentales.

\section{Gestión de Artefactos}

Los resultados se almacenan en una jerarquía estandarizada por conjunto, modo, severidad y método. Cada ejecución preserva métricas agregadas (JSON), señales reconstruidas (NPZ), máscara de daño, configuración inmutable (YAML) y tiempos de ejecución (Tabla~\ref{tab:artifact_types}). El código fuente está versionado con Git, y cada conjunto de resultados está asociado a una etiqueta del repositorio.

\begin{table}[htbp]
\centering
\caption{Artefactos generados por el sistema experimental.}
\label{tab:artifact_types}
\small
\begin{tabular}{@{}l l l@{}}
\toprule
Artefacto & Formato & Función \\
\midrule
Configuración & YAML & Registro inmutable de parámetros \\
Métricas & JSON & Resultados de las seis métricas \\
Señal reconstruida & NPZ & Salida verificable bit a bit \\
Máscara de daño & NPZ & Trazabilidad de canales faltantes \\
Tiempos & JSON & Perfil de eficiencia computacional \\
\bottomrule
\end{tabular}
\end{table}

\begin{table}[htbp]
\centering
\caption{Ejemplos de trazabilidad entre afirmaciones y artefactos.}
\label{tab:traceability}
\small
\begin{tabular}{@{}l l l@{}}
\toprule
Afirmación & Artefacto & Verificación \\
\midrule
TRSS mejora SNR frente a MNE en promedio confirmatorio & Registros de métricas por caso & Diferencia pareada de SNR \\
15K configuraciones & Base de datos Optuna & Conteo de intentos y grillas \\
MNE es más rápido que TRSS & Registros de tiempo & Razón de latencias medianas y medias \\
\bottomrule
\end{tabular}
\end{table}

\section{Decisiones de Diseño}

Tres decisiones arquitectónicas merecen mención explícita. La \textbf{separación estricta de fases} (resultados exploratorios nunca se usaron para ajustar la fase confirmatoria) previene contaminación circular. La \textbf{interfaz común forzada} elimina ventajas espurias por diferencias en preprocesamiento. La \textbf{inmutabilidad de datos originales} garantiza que las diferencias observadas sean atribuibles exclusivamente a los algoritmos de interpolación.